\section{Введение}

В современной обработке естественного языка (NLP) произошел фундаментальный
сдвиг от классических алгебраических алгоритмов к сложным архитектурам на основе
трансформеров. Контекстно-зависимые модели, такие как BERT
(Bidirectional Encoder Representations from Transformers), демонстрируют
высокое качество решения задач семантического поиска, классификации и машинного
понимания текста за счет использования глубоких двунаправленных
представлений. Однако успех этих моделей достигается ценой значительного
усложнения математического аппарата по сравнению с классическими подходами,
такими как латентно-семантический анализ (LSA)

LSA базируется на методах линейной алгебры, в частности на сингулярном
разложении (SVD) терм-документной матрицы. Главным упущением LSA является то,
что семантика фрагмента текста в нем представлена линейной суммой векторов
входящих в него слов. Такой подход полностью игнорирует порядок терминов и
нелинейные синтаксические эффекты, такие как отрицание или логическое следствие.
В результате LSA слова имеют фиксированные векторы, что не позволяет эффективно
решать проблему полисемии (многозначности).

Отличие BERT от LSA заключается в использовании нелинейного математического
аппарата, построенного на механизме самовнимания (self-attention) и полносвязных
сетях с нелинейными функциями активации. Это позволяет модели динамически
вычислять зависимости между словами, формируя контекстуализированные векторы,
которые меняются в зависимости от окружения. Таким образом, BERT не просто
находит наилучшее линейное приближение связей в корпусе, а обрабатывает левый и
правый контексты каждого слова, создавая иерархическую структуру смысла.

Цель настоящей работы - выявить математические основы семантической обработки
текста путем сравнения алгебраических свойств векторных подпространств линейной
(LSA) и нелинейной (BERT) моделей. В статье анализируется, каким образом переход
от линейных матричных преобразований к итеративному механизму внимания позволяет
повысить точность передачи семантики. Практическая значимость работы заключается
в количественной оценке семантической точности моделей, что позволит обосновать
выбор архитектуры в зависимости от требований к учету контекста и доступных
вычислительных мощностей.
